\documentclass{article}

\usepackage{amsmath,amssymb}
\usepackage{xurl}
\usepackage[hidelinks]{hyperref}
\setlength{\emergencystretch}{2em}

% Shared commands for notes in docs/paper-gaps/.

\DeclareUnicodeCharacter{2080}{\ifmmode{}_{0}\else\textsubscript{0}\fi}
\DeclareUnicodeCharacter{2081}{\ifmmode{}_{1}\else\textsubscript{1}\fi}
\DeclareUnicodeCharacter{2082}{\ifmmode{}_{2}\else\textsubscript{2}\fi}
\DeclareUnicodeCharacter{2099}{\ifmmode{}_{n}\else\textsubscript{n}\fi}

\newcommand{\Lean}{\textsc{Lean}}
\ExplSyntaxOn
\NewDocumentCommand{\leanid}{m}{
  \ifmmode
    \textsf{\detokenize{#1}}
  \else
    \group_begin:
    \urlstyle{sf}
    \tl_set:Nn \l_tmpa_tl {#1}
    \tl_replace_all:Nnn \l_tmpa_tl {\_} {_}
    \exp_args:NV \path \l_tmpa_tl
    \group_end:
  \fi
}
\ExplSyntaxOff
\providecommand{\tr}{\operatorname{tr}}
\newcommand{\tnleanIssueBase}{https://github.com/LionSR/TNLean/issues/}
\newcommand{\tnleanPrBase}{https://github.com/LionSR/TNLean/pull/}
\newcommand{\tnleanDocBase}{https://sirui-lu.com/TNLean/docs/}
\newcommand{\ghissue}[1]{\href{\tnleanIssueBase#1}{GitHub issue \##1}}
\newcommand{\ghpr}[1]{\href{\tnleanPrBase#1}{PR \##1}}
\newcommand{\doclink}[1]{\href{\tnleanDocBase#1.html}{\path{#1}}}

% Dirac notation and the identity operator, as in blueprint/src/macros/common.tex.
% \providecommand keeps note-local definitions authoritative.
\usepackage{dsfont}
\providecommand{\ket}[1]{|#1\rangle}
\providecommand{\bra}[1]{\langle #1|}
\providecommand{\braket}[2]{\langle #1 | #2 \rangle}
\providecommand{\ketbra}[2]{|#1\rangle\!\langle #2|}
\providecommand{\Id}{\mathds{1}}
\providecommand{\one}{\mathds{1}}
\providecommand{\C}{\mathbb{C}}
\providecommand{\MN}[1]{M_{#1}(\C)}
\providecommand{\MD}{M_D(\C)}

% Verdict marker: \gapnote{<kind>}{<status>}. Kind is the policy
% classification (clarification, local-correction, scope-restriction,
% unfaithful, false-source, open-gap); status is open, wip (elimination
% underway), resolved, or historical. Machine-read by the site index;
% prints a verdict line.
\newcommand{\gapnote}[2]{\par\noindent\textbf{Verdict:} #1 (#2).\par}


% Fallbacks keep this note compilable when a different command.tex is found first.
\providecommand{\leanid}[1]{\textsf{\detokenize{#1}}}
\providecommand{\gapnote}[2]{\par\noindent\textbf{Verdict:} #1 (#2).\par}

\title{The Cost Model of the Running Time in Landau, Vazirani and Vidick}
\author{}
\date{2026-11-04}

\begin{document}
\maketitle
\gapnote{scope-restriction}{open}

\section{The source statement}

Theorem~1 of Landau, Vazirani and Vidick~\cite{Landau2015Polynomial}
(arXiv:1307.5143v1, \path{algo53.tex}, lines~213--216) asserts that an
algorithm finds a matrix product state within overlap $1-\eta$ of the ground
state ``in time $n^{c(d,\varepsilon)}\,\mathrm{poly}(n/\eta)$''. The paper
does not analyse the linear algebra. Each step is counted as one
polynomial-time computation on data of polynomial dimension: the matrix
elements of the terms of the Hamiltonian in an orthonormal basis of the span
of a family of matrix product states, the solution of a convex program in
polynomially many variables, singular value decompositions, and the leading
eigenvector of a density matrix (lines~219 and 484--487: ``The remaining
steps rely on the singular value decomposition which can be performed
efficiently as well.'').

\section{What is formalized}

The algorithm is written as a program over a type of
queries,\footnote{\leanid{TNLean.MPS.AreaLaw.LinearAlgebraQuery} in
    \doclink{TNLean/MPS/AreaLaw/Algorithm/Query}.} one query for each of
the operations above: one-site extension, enumeration of the net of boundary
contractions, feasibility of the size-trimming convex program and the leading
eigenvector of its solution, Schmidt trimming, the left factors of the
Gaussian filters and their application, and the final convex program. The
honest oracle answers every query with the exact mathematical object that the
definitions of the sweep use, so that the program, run with the honest oracle,
computes the sweep itself and inherits its viability, cardinality and bond
bounds. The cost of a query is declared as one unit plus a polynomial in the
dimensions of its arguments: the chain length, the local dimension, the bond
dimension of the vectors involved, the cardinality of the family, the bond of
the boundary contraction and the rank of the operator. The total declared cost
of the program on a chain of $n$ sites at accuracy parameter $R$ is at most
$P\,(n+2R^2+1)^{d}$ with $P$ and $d$ depending only on the local dimension and
the gap parameter,\footnote{\leanid{TNLean.MPS.AreaLaw.exists_groundStateProgram_cost_polynomial}
    in \doclink{TNLean/MPS/AreaLaw/Algorithm/Cost}.} and the number of
queries is at most the cost.

\section{The restriction}

Three things are declared rather than derived.

\begin{enumerate}
\item \emph{The costs.} A query is charged the cost of the textbook
    algorithm for it on data of the stated dimensions, in exact arithmetic.
    Nothing is proved about how a query is answered. The paper makes the same
    declaration in prose.
\item \emph{The oracle.} The correctness and the cost bound are stated for
    \emph{valid} oracles,\footnote{\leanid{TNLean.MPS.AreaLaw.LinearAlgebraQuery.IsValidOracle},
        \leanid{TNLean.MPS.AreaLaw.groundStateProgram_spec_of_valid} and
        \leanid{TNLean.MPS.AreaLaw.exists_groundStateProgram_cost_polynomial_of_valid}
        in \doclink{TNLean/MPS/AreaLaw/Algorithm/ValidOracle} and
        \doclink{TNLean/MPS/AreaLaw/Algorithm/ValidOracleCost}.} those that
    perform the deterministic operations exactly, enumerate nets of the
    required radius with the standard point bound, report feasibility
    truthfully, return the leading eigenvector of some optimal solution of
    every convex program, return left factors of some cut decomposition of
    every filter with the standard count and rank bounds, and return a unit
    vector of least energy for the final convex program. The honest oracle is
    one of them. The cardinalities of the intermediate families, and with them
    the number of queries, are guaranteed only for valid answers; against an
    arbitrary oracle the enumerated net or the family of filter factors could
    be anything, so a bound over all oracles would require the program to
    truncate every answer to its declared size, which the algorithm does not
    do.
\item \emph{Arithmetic.} Bit complexity and floating-point error are outside
    the model, as they are outside the paper.
\end{enumerate}

The accuracy $\eta$ enters through the padding of the chain by $2R^2$ sites
with a penalty on the padding, in place of the paper's final error-reduction
step at precision $c_\varepsilon/(n\,p(n))$ (lines~669--671): the sweep on the
padded chain at its fixed precision gives a vector of the final span within
distance $1/R$ of the ground state, and the energy-overlap lemma turns the
least energy of the final convex program into overlap at least
$1-4n/(R^2\varepsilon)$. The exponent $c(d,\varepsilon)=2^{O(\log^3 d/\varepsilon)}$
of the source is not reproduced: the library's constant-cut bond is the one
recorded in the note on the constant bond dimension, and the degree $d$ above
is the degree of the uniform polynomial budgets of the sweep.

The padding also changes how the accuracy enters the running time. At the
accuracy parameter $R=\max\{1,\lceil\sqrt{4n/(\varepsilon\eta)}\rceil\}$
the overlap is at least $1-\eta$ and the total declared cost is at most
$P\,(17n/(\varepsilon\eta)+5n)^{d}$ for $\varepsilon\eta\le1$.\footnote{\leanid{TNLean.MPS.AreaLaw.exists_theorem_one_of_valid}
    in \doclink{TNLean/MPS/AreaLaw/Algorithm/TheoremOne}.} This is polynomial
in $n$ and $1/\eta$, and the degree of $1/\eta$ is the degree $d$ of the
whole computation, which depends on the local dimension and the gap. The
printed form $n^{c(d,\varepsilon)}\,\mathrm{poly}(n/\eta)$ could be read as
asserting a degree of $1/\eta$ independent of them, but the source's proof
does not give one: the final error reduction at precision $\varepsilon'/(np(n))$
with $p(n)$ of order $1/\eta$ (lines~669--671) applies the decomposition
(\path{e:decomp}, lines~580--592) with $d^{2\kappa\log(np(n))}\ell$ terms
and $\kappa=O(1/\varepsilon)$, a number of degree of order
$\log d/\varepsilon$ in $1/\eta$. The formalized bound has the same form.

\section{The sampled route}

As printed, the algorithm reduces the error with the sampled approximate
ground-state projector, whose words are drawn uniformly at random, with the
ground energy as an input. The computation on this
route\footnote{\leanid{TNLean.MPS.AreaLaw.sampledGroundStateProgram} and
    \leanid{TNLean.MPS.AreaLaw.exists_sampledGroundStateProgram_success_polynomial}
    in \doclink{TNLean/MPS/AreaLaw/Algorithm/SampledProgram}.} takes one family
of words per cut as an input and is deterministic in them; a sampling
operation is therefore not a query with a declared cost but the law of the
input, the product over the cuts of the law of $\ell$ independent uniform
words. The success probability of Theorem~1 is the probability, under that
law, of the event that every family of words is good, which is at least
$1-1/(2n^2)$ by the union bound over the cuts; on that event the output has
overlap at least $1-\eta$ with the ground state. The size and bond budgets of
this route are polynomial only for a bounded ratio of ground energy to gap,
as recorded in the note on the ground energy in the running time; this
matches the published form of Theorem~1 (Nature Physics 2015, Methods),
whose exponent carries a term of order $\varepsilon_0/\varepsilon$ in
addition to the preprint's $2^{O(\log^3 d/\varepsilon)}$.

\section{Elimination plan}

The declaration of costs is the content of the source's running-time claim
and is not expected to be eliminated. The restriction to the honest oracle
has been lifted on the Gaussian route: the sweep is viable for every optimal
solution of the convex programs and every admissible family of factors, so
the output is correct, with the same cost, for every valid oracle, on the
Gaussian route and on the sampled route alike. 
\bibliographystyle{alpha}
\bibliography{references}

\end{document}
